\documentclass{tutordesk}
\usepackage{td-script}

\studio[NE]{Northbridge English Studio}
\student{Mina Park}
\lesson{7}
\lessondate{26 Oct 2026}
\doctype{Lesson script}
\exam{IELTS Academic \textperiodcentered{} Writing Task 2}
\title{Lesson Script: the PEEL Body Paragraph}
\subtitle{One-to-one, 60 minutes, with the handout of the same name}
\level{Band 6.0 to 7.0}
\timing{60 min}
\focus{Explain move, linking to the question}

\begin{document}
\maketitle

\runningorder

\begin{materials}
  \item Handout \emph{Writing a PEEL Body Paragraph}, two copies (one to write on)
  \item Their last Task 2 essay (libraries question), with one body paragraph marked
  \item Timer; a strip of paper to cover the P/Ex/Eg/L labels on the model
\end{materials}

\begin{teachernote}[Before the lesson]
Their last essay scored about Band 6 for Task Response: ideas are relevant but listed, not
developed. Today has one goal only: every point gets an Explain sentence. Do not correct
grammar today unless meaning breaks down.
\end{teachernote}

% ------------------------------------------------------------------
\stage{5}{Lead-in: what makes an answer convincing?}[Activate the idea that claims need reasons, before any writing.]

\begin{say}
Before we look at essays, a quick question. Your friend says, \enquote{You should move to
Melbourne.} What would they need to say next for you to actually believe them?
\end{say}

\begin{expect}
  \item Why: the weather, jobs, it's cheaper\dots
  \item Give me an example, like a friend who moved there.
\end{expect}

\begin{probe}
  \item So a reason on its own isn't enough. What does the example add?
  \item Which would persuade you more: three reasons, or one reason explained properly?
\end{probe}

\begin{fallback}{If they answer in one word}
Model it yourself in two sentences (claim, then \emph{because}), then ask them to do the same
about their own city.
\end{fallback}

% ------------------------------------------------------------------
\stage{10}{Diagnose: the weak paragraph}[They find the missing moves themselves.]

\begin{action}
  \item Hand out the handout. Cover the model with the paper strip; only the weak paragraph is visible.
  \item Give them two minutes, silently.
\end{action}

\begin{say}
Read the second paragraph, the one on the right. Imagine you are the examiner.
\beat Which sentence actually answers the question about advantages and disadvantages?
\end{say}

\begin{expect}
  \item None of them really.
  \item The last one? \enquote{Libraries should not be replaced\dots}
\end{expect}

\begin{probe}
  \item Does the example about the friend prove anything about cities?
  \item If you deleted the first two sentences, what would the paragraph lose?
\end{probe}

\begin{fallback}{If they defend the weak paragraph}
Agree that the opinion is clear, then ask: \emph{why} should libraries stay? Write their
answer down word for word. That sentence becomes their Explain move in stage 4.
\end{fallback}

% ------------------------------------------------------------------
\stage{12}{Model: the four moves}[Name P, Ex, Eg and L, and see them in a Band 7 paragraph.]

\begin{board}
P \enspace Point \quad one claim that answers the question\par
Ex\enspace Explain \quad why or how it is true\par
Eg\enspace Example \quad one specific case\par
L \enspace Link \quad back to the question's own words
\end{board}

\begin{say}
Now uncover the model on the left. Each sentence has a job. Before you look at my labels,
tell me: which sentence is doing the explaining?
\end{say}

\begin{expect}
  \item The second one, about crowded homes.
\end{expect}

\begin{probe}
  \item What would happen to the example about the teenager if we removed that sentence?
  \item Find the words in the last sentence that come from the question.
\end{probe}

\begin{action}
  \item Read the key point box together. Point out that Explain is the move their essay was missing.
\end{action}

% ------------------------------------------------------------------
\stage{10}{Controlled practice: add the Explain move}[Write one Explain sentence for three given points.]

\begin{say}
I'll give you a point. You give me only the Explain sentence. Use one of the chunks from the
table: \emph{this is largely because} or \emph{which means that}.
\end{say}

\begin{action}
  \item Points, one at a time: (1) Online services save public money. (2) Libraries support
        people without home internet. (3) Digital archives last longer than paper.
  \item Accept the sentence orally first, then they write it.
\end{action}

\begin{expect}
  \item \emph{This is largely because} a single website can serve a whole city, which means
        that councils do not need to pay for buildings and staff.
\end{expect}

\begin{fallback}{If they repeat the point in new words}
Ask \enquote{Why is that true?} and wait. Their spoken answer is the Explain sentence; have them
write exactly what they said.
\end{fallback}

% ------------------------------------------------------------------
\stage{15}{Freer practice: rewrite their own paragraph}[Apply PEEL to the marked paragraph from last week's essay.]

\begin{action}
  \item Give back their essay. Timer: 12 minutes.
  \item While they write, do not interrupt. Note one strong sentence to praise.
\end{action}

\begin{say}
Rewrite this paragraph using all four moves. Keep your original idea; just make the examiner
see your reasoning.
\end{say}

\begin{probe}
  \item Read me only your Explain sentence. Does it say why, or does it add another fact?
  \item Underline the words in your last sentence that come from the question.
\end{probe}

\begin{fallback}{If they finish early}
Ask for a second Example that argues the other side, then a Link that weighs the two.
\end{fallback}

% ------------------------------------------------------------------
\stage{8}{Transfer and wrap-up}[Show that the same structure works in a spoken long turn.]

\begin{cuecard}{Describe a public place in your town that you think is important.}
  \item where it is
  \item who uses it
  \item what people do there
  \andexplain{and explain why you think it is important.}
  \cuenote{One minute to prepare, then speak for up to two minutes.}
\end{cuecard}

\begin{say}
Same idea, but spoken. Give me a point, explain it, give one real example, and finish by
saying why it matters. One minute to think.
\end{say}

\begin{action}
  \item Close with the self-check list on the handout. They tick only what is true of today's paragraph.
  \item Homework: one full Task 2 essay; every body paragraph must have an Explain sentence.
\end{action}

\end{document}
